\documentclass[unicode,11pt,a4paper,oneside,numbers=endperiod,openany, draft]{scrartcl}

% Required package
\usepackage{amssymb}
\usepackage{graphicx}
\usepackage{amsmath}
\usepackage{matlab-prettifier}
\usepackage{float}
\usepackage[export]{adjustbox}
\usepackage{multirow}
\usepackage{booktabs}
\usepackage{amsthm} % math theorems
\usepackage{ifthen}
\usepackage{physics} % responsive norm, abs, ...
\usepackage{algorithm}
\usepackage{algpseudocode}


\renewcommand{\thesubsection}{\arabic{subsection}}

% 1: command name, 2: Title, 3: subtitle, 4: label, 5: content
\newcommand{\mytheorem}[5]{\newtheorem*{#1}{#2} \begin{#1}[#3]\label{#4} #5 \end{#1}}

% 1: if numbered equation, 2: label, 3: content
\newcommand{\myex}[3]{
    \ifthenelse{\equal{#1}{true}}{
        \begin{equation} \label{#2} \begin{aligned} #3 \end{aligned} \end{equation}
    }{
        \begin{equation*} \label{#2} \begin{aligned} #3 \end{aligned} \end{equation*}
    }
}

% vector shortcut
\newcommand{\myvec}[1]{\begin{bmatrix} #1 \end{bmatrix}}

% 1: caption, 2: label, 3: trim, 4: figure within figures folder
\newcommand{\myfigure}[4]{
    \begin{figure}[htbp]
    \centering
    \caption{#1}
    \label{#2}
    \includegraphics[width=\paperwidth, trim=#3]{./figures/#4}
    \end{figure}
}

\usepackage{ifthen}
\usepackage[utf8]{inputenc}
\usepackage{graphics}
\usepackage{graphicx}
\usepackage{hyperref}
\usepackage{amsmath}

\pagestyle{plain}
\voffset -5mm
\oddsidemargin  0mm
\evensidemargin -11mm
\marginparwidth 2cm
\marginparsep 0pt
\topmargin 0mm
\headheight 0pt
\headsep 0pt
\topskip 0pt        
\textheight 255mm
\textwidth 165mm

\newcommand{\duedate} {}
\newcommand{\setduedate}[1]{%
\renewcommand\duedate {\textbf{Due date:}~ #1}}
\newcommand\isassignment {false}
\newcommand{\setassignment}{\renewcommand\isassignment {true}}
\newcommand{\ifassignment}[1]{\ifthenelse{\boolean{\isassignment}}{#1}{}}
\newcommand{\ifnotassignment}[1]{\ifthenelse{\boolean{\isassignment}}{}{#1}}

\newcommand{\assignmentpolicy}{
\begin{table}[h]
\begin{center}
\scalebox{0.8} {%
\begin{tabular}{|p{0.02cm}p{16cm}|}
\hline
&\\
\multicolumn{2}{|c|}{\Large\textbf{Numerical Computing 2023 ---  Submission Instructions}}\\
\multicolumn{2}{|c|}{\large\textbf{(Please, notice that following instructions are mandatory: }}\\
\multicolumn{2}{|c|}{\large\textbf{submissions that don't comply with, won't be considered)}}\\
&\\
\textbullet & Assignments must be submitted to \href{https://www.icorsi.ch/course/view.php?id=14666}{iCorsi} (i.e. in electronic format).\\
\textbullet & Provide both executable package and sources (e.g. C/C++ files, MATLAB). 
If you are using libraries, please add them in the file. Sources must be organized in directories called:\\
\multicolumn{2}{|c|}{\textit{Project\_number\_lastname\_firstname}}\\
& and  the  file must be called:\\
\multicolumn{2}{|c|}{\textit{project\_number\_lastname\_firstname.zip}}\\
\multicolumn{2}{|c|}{\textit{project\_number\_lastname\_firstname.pdf}}\\
\textbullet &  The TAs will grade your project by reviewing your project write-up, and looking at the implementation you attempted, and benchmarking your code's performance.\\

\textbullet & You are allowed to discuss all questions with anyone you like; however: (i) your submission must list anyone you discussed problems with and (ii) you must write up your submission independently.\\
\hline
\end{tabular}
}
\end{center}
\end{table}
}
\newcommand{\punkte}[1]{\hspace{1ex}\emph{\mdseries\hfill(#1~\ifcase#1{Points}\or{Points}\else{Points}\fi)}}


\newcommand\serieheader[6]{
\thispagestyle{empty}%
\begin{flushleft}
\includegraphics[width=0.45\textwidth]{CI_logo}
\end{flushleft}
  \noindent%
  {\large\ignorespaces{\textbf{#1}}\hspace{\fill}\ignorespaces{ \textbf{#2}}}\\ \\%
  {\large\ignorespaces #3 \hspace{\fill}\ignorespaces #4}\\
  \noindent%
  \bigskip
  \hrule\par\bigskip\noindent%
  \bigskip {\ignorespaces {\Large{\textbf{#5}}}
  \hspace{\fill}\ignorespaces \large \ifthenelse{\boolean{\isassignment}}{\duedate}{#6}}
  \hrule\par\bigskip\noindent%  \linebreak
 }

\makeatletter
\def\enumerateMod{\ifnum \@enumdepth >3 \@toodeep\else
      \advance\@enumdepth \@ne
      \edef\@enumctr{enum\romannumeral\the\@enumdepth}\list
      {\csname label\@enumctr\endcsname}{\usecounter
        {\@enumctr}%%%? the following differs from "enumerate"
	\topsep0pt%
	\partopsep0pt%
	\itemsep0pt%
	\def\makelabel##1{\hss\llap{##1}}}\fi}
\let\endenumerateMod =\endlist
\makeatother




\usepackage{textcomp}





\begin{document}

\setassignment
\setduedate{Monday, 3 June 2024, 12:00 AM}

\serieheader
{Optimization Methods}
{2024}
{\textbf{Student:} Jeferson Morales Mariciano \\\\}
{\textbf{Discussed with:} }
{Assignment 4}{}
\newline

%----------------------------------------------------------------------------------
\section{Exercise (20/100)}

Consider the quadratic function \( f: \mathbb{R}^2 \rightarrow \mathbb{R} \) defined as:

\myex{true}{eq:ex1-f}{
    f(\mathbf{x}) = 7 x^2 + 4 xy + y^2
}

\noindent where \( \mathbf{x} = (x, y)^T \).

\begin{enumerate}
    \item Write this function in canonical form, i.e.
          \( f( \mathbf{x} ) = \frac{1}{2} \mathbf{x}^T \mathbf{A} \mathbf{x} + \mathbf{b}^T \mathbf{x} + c \),
          where \( A \) is a symmetric matrix.

    \item Descrive briefly how the Conjugate Gradient (CG) Method works
          and discuss whether it is suitable to minimize \( f \) from equation \ref{eq:ex1-f}.
          Explain your reasoning in detail (max. 30 lines).
\end{enumerate}

\subsection{Answer}

The function written in canonical form correspond to:

\myex{false}{eq:ex1-f-quadratic}{
    f( \mathbf{x} )
    &= 7 x^2 + 4 xy + y^2 \\
    &= \myvec{7x + 2y & 2x + y} \myvec{x \\ y} \\
    &= \myvec{x & y} \myvec{7 & 2 \\ 2 & 1} \myvec{x \\ y} \\
    &= \frac{1}{2} \myvec{x & y} \myvec{14 & 4 \\ 4 & 2} \myvec{x \\ y} \\
    &= \frac{1}{2} \mathbf{x}^T \mathbf{A} \mathbf{x}
}

\noindent With \( \mathbf{b} = \mathbf{0}, \; c = 0 \),
and \( A \) being clearly a symmetric matrix:

\myex{false}{}{
    \mathbf{A} = \myvec{14 & 4 \\ 4 & 2}
}

\noindent Let's verify if \( A \) is positive define as required by the quadratic form:

\myex{false}{}{
    \det \left( \lambda I - \mathbf{A} \right)
    &= \mdet{\lambda - 14 & -4 \\ -4 & \lambda - 2} \\
    &= \lambda^2 - 16 + 12 \\
    &\Rightarrow \lambda_{1,2} = 8 \pm 2 \sqrt{13} > 0
}

\noindent Finally, since all eigenvalues are positive, \( \mathbf{A} \) is SPD.

\subsection{Answer}

The CG method is an algorithm to find the numerical solution of linear equation systems,
with specific application to the minimization of quadratic functions
whose Hessian matrix is symmetric and positive definite.

The CG method is appropriate and effective for minimizing the quadratic function.


%----------------------------------------------------------------------------------
\section{Exercise (20/100)}

Consider the following constrained minimization problem for \( \mathbf{x} = (x, y, z)^T \)

\myex{true}{eq:ex2-min-f}{
    &\min_{\mathbf{x}} f(\mathbf{x}) := -3x^2 + y^2 + 2z^2 + 2(x + y + z) \\
    &\text{subject to} \quad c(\mathbf{x}) = x^2 + y^2 + z^2 - 1 = 0
}

\noindent Write down the Lagrangian function and derive th KKT conditions for \ref{eq:ex2-min-f}




%----------------------------------------------------------------------------------
\section{Exercise (60/100)}

\begin{enumerate}
    \item Read the chapter on Simplex method, in particular the section 13.3 The Simplex Method,
          in Numerical Optimization, Nocedal and Wright.
          Explain how the method works, with a particular attention to the search direction.

    \item Consider the following contrained minimization problem,
          \( \mathbf{x} = (x_1, x_2)^T \);

          \myex{true}{eq:ex3-min-f}{
              \min_{\mathbf{x}} f(\mathbf{x}) := 4x_1 + 3x_2
          }

          subect to:

          \myex{true}{eq:ex3-constraints}{
              6 - 2x_1 - 3x_2 \geq 0 \\
              3 + 3x_1 - 2x_2 \geq 0 \\
              5 - 2x_2 \geq 0 \\
              4 - 2x_1 - x_2 \geq 0 \\
              x_2 \geq 0 \\
              x_1 \geq 0 \\
          }

          \begin{enumerate}
              \item Sketch the feasible region for this problem.

              \item Which are the basic feasible points of the problem \ref{eq:ex3-min-f}?
                    Compute them by hand using the geometrical interpretation
                    and find the optimal point \( \mathbf{x^*} \)
                    that minimizes \( f \) subject to the constraints.

              \item Prove that the first order necessary conditions holds for the optimal point.
          \end{enumerate}
\end{enumerate}

\end{document}
